% Figura: polinomios de Legendre y serie no truncada para l = 1
\begin{figure}[ht]
\centering
\begin{tikzpicture}
  \begin{groupplot}[group style={group size=2 by 1, horizontal sep=1.3cm},
      oscuro, width=0.5\textwidth, height=6.2cm, xmin=-1, xmax=1,
      xtick={-1,-0.5,0,0.5,1}, xlabel={$x$}]
    % ---------- (a) P_0 ... P_4 ----------
    \nextgroupplot[ymin=-1.1, ymax=1.1, ytick={-1,-0.5,0,0.5,1},
        title={(a) Polinomios $P_\ell$ (series truncadas)},
        samples=200, domain=-1:1,
        legend style={at={(0.5,-0.24)}, anchor=north, legend columns=5,
                      /tikz/every even column/.append style={column sep=3pt}}]
      \addplot[line width=1.15pt, gris] {1};
      \addplot[c1] {x};
      \addplot[c2] {0.5*(3*x^2-1)};
      \addplot[c3] {0.5*(5*x^3-3*x)};
      \addplot[c4] {(35*x^4-30*x^2+3)/8};
      \legend{$P_0$, $P_1$, $P_2$, $P_3$, $P_4$}
    % ---------- (b) y1 = 1 - x artanh x ----------
    \nextgroupplot[ymin=-2.6, ymax=1.25, ytick={-2,-1,0,1},
        title={(b) Serie par para $\ell=1$: no se corta},
        legend style={at={(0.5,-0.24)}, anchor=north, legend columns=5,
                      /tikz/every even column/.append style={column sep=3pt}}]
      \draw[guia] (axis cs:-1,-2.6) -- (axis cs:-1,1.25);
      \draw[guia] (axis cs:1,-2.6) -- (axis cs:1,1.25);
      \addplot[c2] table[x=x, y=S1] {datos/legendre_l1_parciales.dat};
      \addplot[c3] table[x=x, y=S2] {datos/legendre_l1_parciales.dat};
      \addplot[c4] table[x=x, y=S4] {datos/legendre_l1_parciales.dat};
      \addplot[c5] table[x=x, y=S8] {datos/legendre_l1_parciales.dat};
      \addplot[blanco, line width=1.6pt, samples=400, domain=-0.997:0.997]
        {1-0.5*x*ln((1+x)/(1-x))};
      \legend{$S_1$, $S_2$, $S_4$, $S_8$, $y_1$}
  \end{groupplot}
\end{tikzpicture}
\caption{(a) Para $\ell\in\N$ la serie de la paridad de $\ell$ se corta y da
el polinomio de Legendre $P_\ell$, regular en todo $[-1,1]$ (y, de hecho, en
todo $\R$). (b) Para $\ell=1$, la serie par
$y_1=1-x\artanh x$ de \eqref{eq:legendre-q1} y sus sumas parciales
$S_M=1-\sum_{k=1}^{M}x^{2k}/(2k-1)$: la convergencia se degrada al
acercarse a $x=\pm1$, donde $y_1$ diverge logarítmicamente. El radio $R=1$
del \cref{teo:fuchs} se alcanza.}
\label{fig:legendre}
\end{figure}
